\documentclass[12pt,a4paper]{article}
\usepackage[utf8]{inputenc}
\usepackage[T1]{fontenc}
\usepackage{lmodern}
\usepackage[margin=2.5cm]{geometry}
\usepackage{amsmath,amssymb}
\usepackage{booktabs}
\usepackage{hyperref}

\title{CCNN v0.8: Candidate Loss ve Model Loss (NaN/Inf) Koken Analizi}
\author{GitHub Copilot (GPT-5.3-Codex)}
\date{7 Nisan 2026}

\begin{document}
\maketitle

\section{Amac}
Bu rapor, su soruyu netlestirmek icin hazirlanmistir:
\begin{quote}
Hidden unit (ozellikle 2. hidden) eklendikten sonra gorulen model loss NaN/Inf problemi,
candidate tarafindaki buyuk loss ile nasil iliskilidir?
\end{quote}

\section{Incelenen Akis}
Egitim akisinin ilgili adimlari:
\begin{enumerate}
    \item Candidate egitimi yapilir (candidate objective).
    \item En iyi candidate secilir ve aga eklenir.
    \item Output layer, yeni hidden ile tekrar MSE loss uzerinden egitilir.
    \item Sonra recursive full-series MSE hesaplanir.
\end{enumerate}

Kodda bu akis, \texttt{Npred\_MiniBatch\_Adam\_maxCandidate.m} icinde gorulur:
\begin{itemize}
    \item Candidate train cagrisi
    \item \texttt{W\_hidden\{end+1\} = w\_h} ile hidden ekleme
    \item \texttt{trainOutputLayer\_Trajectory(...)} ile output retrain
    \item \texttt{current\_mse = mean((Ytr(2:end)-Yhat\_tmp(2:end)).\^{}2)}
\end{itemize}

\section{Loss Tanimlari}
\subsection{Candidate tarafi}
Candidate loss su sekilde tanimli:
\begin{equation}
L_{cand} = -(\text{metric})^2
\end{equation}

Metric ise su sekilde kullaniliyor:
\begin{equation}
\text{metric} = |\text{cov}_{vr}|, \quad
\text{cov}_{vr} = \sum_i v_{c,i} r_{c,i}
\end{equation}

Burada onemli nokta: \( |\text{cov}_{vr}| \) sinirsizdir (unbounded).
Yani veri olcegi, residual olcegi, veya dinamik davranis ile buyuyebilir.

\subsection{Model (output) tarafi}
Output retrain asamasinda optimize edilen loss:
\begin{equation}
L_{out} = \text{MSE}(Y, T)
\end{equation}

Bu loss, forwardModelTrajectory ile uretilen kapali-cevrim tahminlerden gelir.

\section{Temel Bulgular}
\subsection{Dogrudan mi, dolayli mi?}
\textbf{Dogrudan degil, dolayli.}

Candidate loss un buyumesi tek basina model loss u NaN yapmaz.
Ancak buyuk ve agresif candidate secimi, eklendikten sonra kapali-cevrim ileri geciste
sayisal kararsizlik olusturup MSE tarafini NaN/Inf e goturebilir.

\subsection{Neden hidden eklendikten sonra patliyor?}
Hidden eklendikten sonra su zincir calisir:
\begin{align}
\text{(1) Candidate secimi} &\Rightarrow \text{buyuk etkili } w_h \\
\text{(2) Hidden ekleme} &\Rightarrow x \mapsto [x, a_h] \\
\text{(3) Output retrain + feedback} &\Rightarrow y_t \text{ tekrar } y\text{-history icine girer} \\
\text{(4) Amplifikasyon} &\Rightarrow y_t \to \pm\infty \text{ veya NaN} \\
\text{(5) MSE} &\Rightarrow \text{NaN/Inf}
\end{align}

Yani NaN gordugunuz yer model tarafidir; koken etkisi candidate secim kriterinden gelebilir.

\subsection{Clipping neden tek basina cozum degil?}
Gradient clipping (varsa) update adimini sinirlar,
ama objective fonksiyonunun olcegini sinirlamaz.
Unbounded metric hedeflenmeye devam ettigi icin, yine kararsiz bir candidate secilebilir.

\section{Sizin Sorunuza Net Cevap}
\begin{quote}
\textit{``Ben hidden i aga ekleyip sonra tekrar loss hesaplarken NaN goruyorum.''}
\end{quote}

Bu gozlem teknik olarak dogrudur.
NaN gorunen yer output/model MSE hattidir.
Ancak bu NaN, onceki candidate secim asamasinda secilen hidden in
kapali-cevrim davranisi bozmasi nedeniyle tetiklenebilir.

Kisa ifade ile:
\begin{equation}
\text{Candidate objective mismatch} \;\Longrightarrow\; \text{agresif hidden secimi}
\;\Longrightarrow\; \text{output retrain/forward pass NaN}
\end{equation}

\section{Sonuc}
Problem sadece ``candidate loss buyuk'' veya sadece ``model loss NaN'' degildir;
ikisi bir nedensel zincirin farkli halkalaridir.

\begin{itemize}
    \item NaN nin gorundugu yer: model/output MSE hatti.
    \item NaN yi tetikleyebilen koken: candidate secim hedefinin MSE ile tam hizali olmamasi ve unbounded metric kullanimi.
\end{itemize}

\end{document}
